\begin{thebibliography}{99}
%======================================================================
\bibitem{fiedler1973} M.~Fiedler. Algebraic connectivity of graphs.
\emph{Czechoslovak Mathematical Journal}, 23(2):298--305, 1973.

\bibitem{ghosh2006} A.~Ghosh and S.~Boyd. Growing well-connected graphs. In
\emph{Proc.\ 45th IEEE Conference on Decision and Control (CDC)}, 6605--6611, 2006.

\bibitem{pecora1998} L.~M.~Pecora and T.~L.~Carroll. Master stability functions
for synchronized coupled systems. \emph{Physical Review Letters},
80(10):2109--2112, 1998.

\bibitem{barahona2002} M.~Barahona and L.~M.~Pecora. Synchronization in
small-world systems. \emph{Physical Review Letters}, 89(5):054101, 2002.

\bibitem{bomela2020} W.~Bomela, S.~Wang, C.-A.~Chou, and J.-S.~Li. Real-time
inference and detection of disruptive EEG networks for epileptic seizures.
\emph{Scientific Reports}, 10:8653, 2020.

\bibitem{weyl1912} H.~Weyl. Das asymptotische Verteilungsgesetz der Eigenwerte
linearer partieller Differentialgleichungen (mit einer Anwendung auf die Theorie
der Hohlraumstrahlung). \emph{Mathematische Annalen},
71(4):441--479, 1912.

\bibitem{daviskahan1970} C.~Davis and W.~M.~Kahan. The rotation of eigenvectors
by a perturbation.\ III. \emph{SIAM Journal on Numerical Analysis}, 7(1):1--46, 1970.

\bibitem{delong1988} E.~R.~DeLong, D.~M.~DeLong, and D.~L.~Clarke-Pearson.
Comparing the areas under two or more correlated receiver operating
characteristic curves: a nonparametric approach. \emph{Biometrics},
44(3):837--845, 1988.

\bibitem{sunxu2014} X.~Sun and W.~Xu. Fast implementation of DeLong's algorithm
for comparing the areas under correlated receiver operating characteristic
curves. \emph{IEEE Signal
Processing Letters}, 21(11):1389--1393, 2014.

\bibitem{mitchell2003} C.~C.~Mitchell and D.~G.~Schaeffer. A two-current model
for the dynamics of cardiac membrane. \emph{Bulletin of Mathematical Biology},
65(5):767--793, 2003.

\bibitem{roney2022} C.~H.~Roney et al. Predicting atrial fibrillation recurrence
by combining population data and virtual cohorts of patient-specific left atrial
models. \emph{Circulation: Arrhythmia and Electrophysiology}, 15(2):e010253, 2022.
doi:10.1161/CIRCEP.121.010253. Data: Zenodo 5801337.

\bibitem{uwboyle2025} S.~F.~Bifulco, M.~J.~Magoon, Y.~Chahine, I.~Kim,
F.~Macheret, N.~Akoum, and P.~M.~Boyle. Predicting arrhythmia
recurrence post-ablation in atrial fibrillation using explainable machine
learning. \emph{Communications Medicine}, 5(1):421, 2025.
doi:10.1038/s43856-025-01058-4. PMID 41087542.
Data: Dryad \code{10.5061/dryad.kkwh70sg0}.

\bibitem{khurram2014} I.~M.~Khurram et al. Magnetic resonance image intensity
ratio, a normalized measure to enable interpatient comparability of left atrial
fibrosis. \emph{Heart Rhythm}, 11(1):85--92, 2014.

\bibitem{kumar2012} S.~Kumar et al. Atrial fibrillation inducibility in the
absence of structural heart disease or clinical atrial fibrillation: critical
dependence on induction protocol, inducibility definition, and number of
inductions. \emph{Circulation: Arrhythmia and Electrophysiology}, 5(3):531--536,
2012. PMID 22528047.

\bibitem{marquardt2018} J.~Marquardt et al. Prognostic value of atrial
fibrillation inducibility in patients without history of clinical atrial
fibrillation. \emph{Journal of Atrial Fibrillation}, 11(1):1837, 2018.
PMID 30455834.

\bibitem{oral2008} H.~Oral et al. Inducibility of paroxysmal atrial fibrillation
by isoproterenol and its relation to the mode of onset of atrial fibrillation.
\emph{Journal of Cardiovascular Electrophysiology}, 19(5):466--470, 2008.
PMID 18266669.

\bibitem{kawai2019} S.~Kawai et al. Predictive value of the induction test with
atrial burst pacing with regard to long-term recurrence after ablation in
persistent atrial fibrillation. \emph{Journal of Arrhythmia}, 35(2):223--229,
2019. PMID 31007786.

\bibitem{liu2020} H.~Liu et al. Is atrial fibrillation noninducibility by burst
pacing after catheter ablation associated with reduced clinical recurrence? A
systematic review and meta-analysis. \emph{Journal of the American Heart
Association}, 9(14):e015260, 2020. PMID 32654581.

\bibitem{tripodai2024} G.~S.~Collins, K.~G.~M.~Moons, P.~Dhiman et al.
TRIPOD+AI statement: updated guidance for reporting clinical prediction models
that use regression or machine learning methods. \emph{BMJ}, 385:e078378, 2024.
PMID 38626948.
\end{thebibliography}